\begin{proposition}\label{prop:decic-no93}
An elliptic curve over a number field of degree $10$ has no point
of order $93$.
\end{proposition}
\begin{proof}
Put $X=X_1(93)$ and $J=J_1(93)$, of genus $261$.
The exact winding decomposition has one winding-zero constituent,
of dimension $2$ and trivial character. Every other constituent has
nonzero winding. By the winding rank-zero criterion of \cite[Section~3]{DEHMZ},
$(g-1)J(\mathbb Q)$ is finite for
every diamond $g$. No rank-zero assertion for $J$ is needed.
Let $D$ be the trace divisor of a hypothetical level structure over
a field of degree $10$. Its source point has exact degree $10$:
the proper divisors $1,2,5$ of $10$ are excluded by the absence of
order $31$ in those degrees. Thus $D$ is reduced.

The independent complete finite-field census gives no noncuspidal
places in degrees at most $8$, ten in degree $9$, and none in degree
$10$. The degree-$9$ places form one diamond orbit. Their curves have
$558$ points, cyclic group, and one closed $j$-orbit of degree $3$.
Hence noncuspidal reduction is a degree-$9$ place together with a
rational cusp. Put $\sigma=\langle2\rangle^{-1}$ and
$\tau=\langle5\rangle$, using point levels.
The finite diamond image and Manin--Drinfeld give rational torsion
$[\sigma D-D]$, and the rational cuspidal complement also gives
rational torsion. The operator $A_5=T_5-1-5\tau$ annihilates it.
Thus, for $P$ the noncuspidal residual place and $T=T_5P$,
\[
 B=\sigma T-T-\sigma P+P-5\sigma\tau P+5\tau P
\]
must be principal. The accepted exact canonical multiplication
certificate proves it nonprincipal for the orbit representative and
all necessary diamond translates. It uses the actual six geometric
$5$-kernels, with three degree-$18$ target places and norm
multiplicity $1$. The signed parts have degree $108$, and every
necessary canonical multiplication matrix has rank $261$.
This obstruction depends only on $P$, and remains
valid with the present rational cuspidal complement. Hence we are
left with entirely cuspidal reduction.

The finite cusp signature is $1^{30}2^{15}5^6 10^3$.
We first dispose of the quadratic and decic cusps. Put
$h=\langle32\rangle$, an involution. Exact Tate-label arithmetic
shows that $h$ fixes every quadratic and every decic closed cusp.
It interchanges pairs of quintic cusps.
Lift $\overline D$ to a cuspidal divisor $C$ over $\mathbb Q_2$.
The cusps are unramified at $2$; the lift need not be rational over
$\mathbb Q$. The class $(h-1)[D-C]$ is torsion, by finite rational
image and Manin--Drinfeld, and reduces to zero. The torsion reduction
kernel over $\mathbb Q_2$ is killed by $2$, so
\[
 2hD+2C\sim2D+2hC.
\]
If $C$ contains a quadratic or decic closed cusp, then
$\deg\gcd(C,hC)\geq2$. After cancellation, the displayed divisors
have degree at most $36$. However, the projective index is $3840$,
and \cite[(2.1)]{NV} gives
\[
 \operatorname{gon}_{\mathbb C}X>
 \frac{325\cdot3840}{32768}=\frac{4875}{128}>36.
\]
The two divisors are equal. Their noncuspidal parts imply $hD=D$.

The involution $h$ is free on geometric noncuspidal points in
characteristic zero. For generic $j$ the projective stabilizer is
trivial. At $j=0$ it has order $3$, so cannot contain the projective
class of $h$, which has order $2$. At $j=1728$ a nontrivial projective
stabilizer would be induced by an automorphism whose square is $-1$.
This is incompatible with $32^2\equiv1\pmod{93}$ on a point of odd
order $93$.
Since $D$ is a single degree-$10$ closed trace, $h$ therefore induces
an involution on its residue field. The fixed field has degree $5$.
The forgetful map $(E,P)\mapsto(E,3P)$ to $X_1(31)$ is $h$-invariant,
since $32\equiv1\pmod{31}$. Thus a point of order $31$ descends to
a field of degree $1$ or $5$, contrary to the quintic torsion
classification. This excludes every reduction containing a
quadratic or decic cusp.

It remains to treat rational and quintic cusps. We use the
multiplicative-level canonical model. Its diamonds for $2,5$ have
finite rational difference image. Let
\[
 \Psi_2=(T_2^\mu-1)(\delta_2-1),\qquad
 \Psi_5=(T_2^\mu-1)(\delta_5-1).
\]
For a cuspidal lift $C$ as above, each $(\delta_i-1)[D-C]$ is torsion
over $\mathbb Q_2$ and has zero reduction. The multiplicative-level
version of the full-Jacobian dyadic kernel lemma gives
$\Psi_i[D-C]=0$. The factor is $T_2^\mu-1$, obtained from the
point-level factor $T_2-\langle2\rangle$ by Fricke and an invertible
diamond. All cotangent calculations take place on the full Jacobian.

The pullback image of these two operators is a diamond-stable space
$U$ of dimension $248$ over $\mathbb F_2$. Translate a rational cusp
of greatest multiplicity to the first cusp and retain ten jets there
and five at the other $29$ rational cusps. The $155$ rows have rank
$150$, and their full relation space has dimension $5$. Independent
enumeration of all $31$ nonzero relations gives minimum normalized
prefix cost $30$. Thus every rational-cuspidal divisor of degree
$10$ gives restriction rank $10$.

The six quintic cusps form one diamond orbit. A genuine Tate family
for one of them is
\[
                 E(q^3),\qquad P=\zeta_{31}q,
\]
over $\mathbb F_{32}$. The accepted global modular-unit identities
recover its canonical arc through the first two coefficients.
Two different full-rank minors give the same arc; the tangent is
nonzero. The diamond action is identified globally with inverse
point multiplication. Coefficientwise Frobenius on this arc is
checked against point multiplication by $64$, which is $1$ modulo
$3$ and $2$ modulo $31$. This verifies the six closed-cusp restriction
blocks, including their residue-field descent.

If a quintic cusp occurs, move one to the first quintic place.
The remaining possibilities are one quintic plus rational degree
$5$, or two quintics, including a double cusp. We test
\[
                    \binom{34}{5}+6=278262
\]
normalized configurations. Every matrix has rank $10$.
A separate complete-matrix implementation verifies these ranks and
the rational relation space. The canonical unit identities used for
these arcs are global: their verified orders exceed the section
degree bounds coming from unit pole degrees $100,105$.

Thus the product of $\Psi_2,\Psi_5$ is a formal immersion on
$X^{(10)}$ at every remaining cuspidal reduction. Since $D$ and $C$
have the same reduction and the same image, we obtain $D=C$.
This contradicts the noncuspidal support of $D$ and completes the
proof.
\end{proof}
